% ============================================================================
% Section 6.6 — Angle Relationships: Supplementary, Complementary, and
%                Vertical Angles
% CCSS 7.G.B.5
% ============================================================================
\section{Angle Relationships}

\begin{stepGoal}
\begin{itemize}
    \item Identify complementary, supplementary, and vertical angle pairs
    \item Write and solve equations to find unknown angles
\end{itemize}
\end{stepGoal}

\begin{stepVocabBox}
    \vocabItem{Complementary}{Two angles that add to $90°$.}
    \vocabItem{Supplementary}{Two angles that add to $180°$.}
    \vocabItem{Vertical Angles}{Opposite angles formed by two intersecting lines --- always equal.}
\end{stepVocabBox}

\begin{stepsCard}[How to Find an Unknown Angle]
    \stepItem{Identify the \textbf{angle relationship}: complementary ($90°$), supplementary ($180°$), or vertical (equal).}
    \stepItem{Write an \textbf{equation} using the relationship.}
    \stepItem{\textbf{Solve} for the variable and find each angle measure.}
    \stepItem{\textbf{Check}: do the angles satisfy the original relationship?}
\end{stepsCard}

% --- Worked Example 1 ---
\begin{stepExample}{Two supplementary angles are $3x + 10$ and $2x$. Find both angles.}
    \stepShow{1} Supplementary $\to$ sum is $180°$.

    \stepShow{2} Equation: $3x + 10 + 2x = 180$.

    \stepShow{3} Combine: $5x + 10 = 180 \to 5x = 170 \to x = 34$.\par
    Angles: $3(34) + 10 = 112°$ and $2(34) = 68°$.

    \stepShow{4} Check: $112 + 68 = 180°$ \checkmark
    \stepResult{$112°$ and $68°$}
\end{stepExample}

% --- Worked Example 2 ---
\begin{stepExample}{Two vertical angles are $4x + 5$ and $6x - 15$. Find $x$ and the angle.}
    \stepShow{1} Vertical angles are \textbf{equal}.

    \stepShow{2} Equation: $4x + 5 = 6x - 15$.

    \stepShow{3} Subtract $4x$: $5 = 2x - 15$. Add $15$: $20 = 2x$. So $x = 10$.\par
    Angle: $4(10) + 5 = 45°$.

    \stepShow{4} Check: $6(10) - 15 = 45°$ \checkmark
    \stepResult{$x = 10$; both angles are $45°$.}
\end{stepExample}

\watchOut{Don't mix up complementary ($90°$) and supplementary ($180°$). \textbf{C} comes before \textbf{S} in the alphabet, and $90 < 180$.}

% ============================================================================
% PRACTICE
% ============================================================================
\newpage
\begin{stepPractice}{Angle Relationships Practice}
    \resetProblems

    \stepPracticeHeader[step1Color]{Name the Relationship}
    \begin{multicols}{2}
    \prob Two angles sum to $90°$. They are \answerBlank[2.5cm].
    \answer{complementary}
    \prob Two opposite angles from intersecting lines are \answerBlank[2.5cm].
    \answer{vertical (equal)}
    \end{multicols}

    \stepPracticeHeader[step2Color]{Write and Solve Equations}
    \prob Two complementary angles: one is $35°$. Find the other. \answerBlank[2cm]
    \answer{$55°$}

    \prob Two supplementary angles: one is $128°$. Find the other. \answerBlank[2cm]
    \answer{$52°$}

    \stepPracticeHeader[step3Color]{Solve with Variables}
    \prob Supplementary angles $(x + 30)°$ and $(2x)°$. Find both. \answerBlank[3cm]
    \answer{$80°$ and $100°$}

    \wordProblem{An angle and its complement are in the ratio $2{:}3$. Find both angles.}{~}
    \answer{$36°$ and $54°$. Let the angles be $2x$ and $3x$: $5x = 90$, $x = 18$.}
\end{stepPractice}
